\documentclass[10pt,a4paper]{amsart}
\usepackage[margin=19mm]{geometry}
\usepackage{amsmath,amssymb,mathtools}
\usepackage[scr=rsfs,cal=euler]{mathalfa}
\usepackage{xfrac}
\usepackage[unicode,hidelinks,bookmarks=false]{hyperref}
\newcommand{\ie}{i.e.\ }
\newcommand{\cf}{cf.\ }
\newcommand{\resp}{respectively\ }
\newcommand{\Subsection}{\subsection}
\providecommand{\nobreakdash}{\nobreak\hbox{-}}
\setlength{\emergencystretch}{2em}
\multlinegap0pt
\usepackage{amssymb}
\usepackage{xcolor}
\newtheorem{lemma}{Lemma}
\newtheorem{theorem}{Theorem}
\DeclareMathOperator{\Perm}{Perm}
\title{Permutohedral complex and complements of diagonal subspace arrangements}
\author{Taras Panov, Vsevolod Tril}
\date{Source: arXiv:2504.04183v2 (CC BY 4.0)}
\begin{document}
\maketitle

\noindent\emph{Source and licence.} Permutohedral complex and complements of diagonal subspace arrangements. Version arXiv:2504.04183v2 (CC BY 4.0), distributed by arXiv under the Creative Commons Attribution 4.0 International licence, http://creativecommons.org/licenses/by/4.0/, as recorded in the arXiv licence metadata for this exact version and shown on the abstract page https://arxiv.org/abs/2504.04183v2. Full licence notice: \texttt{licenses/arxiv-2504.04183-CC-BY-4.0.txt}.\par\smallskip

\noindent\textit{Editorial scope.} Counted unit: the article's theorem theorem (source0.tex lines 801--804). The article's complete original proof of it is delivered with it (source0.tex lines 805--845, one \texttt{proof} environment of 41 lines). No mathematics is written, repaired or re-derived: the statement and the proof are the article's own lines, in the article's own notation, and no retained line is substituted or re-flowed.  Retained uncounted as the source-local context the proof needs: lines 260--262; lines 680--682; lines 684--687; lines 726--728; lines 767--773.  Nothing was omitted from any retained span, so the package carries no omission record.  No retained line contains a citation, so the wrapper carries no bibliography and no bibliography entry.  No retained line contains a figure environment, an image-inclusion command or drawing code, so no image is delivered and the package declares no asset dependency.  Editorial formatting only, in addition: an AMS \texttt{amsart} preamble that restates the article's own declarations and macros used by the retained lines, namely the environments \texttt{lemma}, \texttt{theorem} on separate counters, together with the standard packages those lines need.  The retained environments print in this document as Lemma 1, Theorem 1 in order of appearance, whereas the article prints the same items with the numbering of its own full document.  The complete native source of the article as distributed in the arXiv e-print for this exact version is bundled unchanged as \texttt{sources/176040/source0.tex}.
\begin{equation}\label{DeltaC}
  \Delta_{C}(u_{[m]} t_\varnothing \underline{t}_{\varnothing}) = \sum_{\sigma \subset [m]} (-1)^{(\sigma, [m] \backslash \sigma)} u_{\sigma} t_\varnothing \underline{t}_{[m] \backslash \sigma} \otimes u_{[m] \backslash \sigma} t_{\sigma} \underline{t}_{\varnothing}.
\end{equation}  

\begin{equation}\label{SU-diag}
\Delta_{SU}(F(1, 2, \ldots, m)) = \sum_{q = 1}^m \sum_{A \in \mathcal{C}^{q \times (m-q+1)}} \mathop{\mathrm{csgn}}(A)\; F(c(A)) \otimes F(r(A))
\end{equation}

\begin{equation}\label{comulrule}
  \Delta_{SU}\bigl(F(U_1| \cdots | U_p)\bigr) = 
  F\bigl(\Delta_{SU}(F(U_1)) \big| \cdots \big| \Delta_{SU}(F(U_p))\bigl).
\end{equation}

\begin{equation}\label{rho-face}
\rho \left({F(U_1 | \cdots | U_p) }\right) = \prod_{i \in \sigma}D^1 \times \prod_{i \in \tau} \{ 1 \} \times \prod_{i \notin \sigma \cup \tau} \{-1\},
\end{equation}

\begin{lemma}\label{snake}
Suppose that a matrix $A \in \mathcal{O}^{q \times (m-q+1)}$ 
is such that the dimension of the faces 
$F(c(A))$ and $F(r(A))$ is preserved under the projection $\rho$.
Then all nonzero entries of $A$ form one continuous snake. 

\end{lemma}

\begin{theorem}
For any face $F(U_1 | \cdots | U_p)$ of $\Perm^{m-1}$ we have
$$(\rho_* \otimes \rho_*) \Delta_{SU} F(U_1 | \cdots | U_p) = \Delta_{C}(\rho_* F(U_1 | \cdots | U_p)).$$
\end{theorem}

\begin{proof}
By the comultiplicative rule~\eqref{comulrule} it is sufficient to consider the action of the homomorphism on the face of maximal dimension.
Consider $\Delta_{SU}F([m])$. By Lemma~\ref{snake}, the image of $\rho_* \colon C_{m-1}(\Perm^{m-1}) \to C_{m-1}(I^{m-1})$ consists of those
summands in~\eqref{SU-diag} that correspond to configuration matrices with one continuous snake.

We claim that all matrices satisfying the hypothesis of Lemma~\ref{snake} are configuration matrices.
Indeed, every ordered matrix $A$ containing one continuous snake with nodes $1=i_0, i_1, \ldots, i_r = m$ can be derived from a step matrix $E$ of the following form. 
Suppose that the elements $\{1, 2, \ldots, i_1 \}$ form a row and $r$ is even.
The first row of $E$ is formed by entries $\{1, 2, \ldots, i_1, i_2+1, i_2 + 2, \ldots, i_3, i_4+1, i_4 + 2, \ldots, i_{r - 1}\}$ arranged in consecutive columns, and the first column of $E$ is formed by entries
$\{1, i_1 + 1, i_1 + 2, \ldots, i_2, i_3 + 1, i_3 + 2, \ldots, i_r \}$ arranged in consecutive rows. (If $r$ is odd, we should swap $i_{r-1}$ and $i_r$, and if entries $\{1, 2, \ldots, i_1 \}$ form a column, we should transpose the matrix.)

At the $j$th step, we will shift the column of elements $\{i_j + 1, \ldots, i_r \}$ to the right until we reach the column containing the element $i_j$. Then at the $(j+1)$th step, we shift the row of elements $\{i_{j+1} + 1, \ldots, i_{r-1} \}$ down until the row containing element the $i_{j+1}$. The procedure is then performed until we obtain the matrix $A$. The first two steps of this procedure are as follows:
$$\begin{pmatrix}
1 & 2 & \cdots & i_1 & i_2 + 1 & \cdots & i_{r-1} \\
i_1 + 1 & 0 & \cdots & 0 & 0 & \cdots & 0 \\
i_1 + 2 & 0 & \cdots & 0 & 0 & \cdots & 0 \\
\vdots & \vdots & \cdots & \vdots & \vdots & \cdots & \vdots\\
i_2 & 0 & \cdots & 0 & 0 & \cdots & 0 \\
\vdots & \vdots & \cdots & \vdots & \vdots & \cdots & \vdots \\
i_{r} & 0 & \cdots & 0 & 0 & \cdots & 0 
\end{pmatrix} \to
\begin{pmatrix}
1 & 2 & \cdots & i_1 & 0 & \cdots & 0 \\
0 & 0 & \cdots & i_1 + 1 & 0 & \cdots & 0 \\
0 & 0 & \cdots & i_1 + 2 & 0 & \cdots & 0 \\
\vdots & \vdots & \cdots & \vdots & \vdots & \cdots & \vdots\\
0 & 0 & \cdots & i_2 & i_2 + 1 & \cdots & i_{r-1} \\
\vdots & \vdots & \cdots & \vdots & \vdots & \cdots & \vdots \\
0 & 0 & \cdots & i_{r} & 0 & \cdots & 0 
\end{pmatrix}$$

Let $\sigma = \{i_1, \ldots, i_2 - 1, i_3, \ldots, i_4 - 1, \ldots, i_{r} - 1\}$ and $\sigma' = \{1, \ldots, i_1-1, i_2, \ldots, i_3-1, \ldots, i_{r} \}$. Relation~\eqref{rho-face} implies that $\rho(F(c(A))) = u_{\sigma} t_\varnothing \underline{t}_{\sigma'}$ and
$\rho(F(r(A))) = u_{\sigma'} t_{\sigma} \underline{t}_\varnothing$.
Since each subset $\sigma \subset [m]$ corresponds to a single snake, we have
\[
  (\rho_* \otimes \rho_*)\Delta_{SU}F([m]) = \sum_{\sigma \subset [m]} 
  (-1)^\varepsilon u_\sigma t_\varnothing \underline{t}_{\sigma'} \otimes 
  u_{\sigma'}t_{\sigma} \underline{t}_\varnothing.
\]
Direct calculations show that the sign $(-1)^\varepsilon$ coincides with the sign in formula~\eqref{DeltaC} for~$\Delta_{C}$.
\end{proof}

\end{document}
